% Part: computability
% Chapter: computability-theory
% Section: k-1

\documentclass[../../../include/open-logic-section]{subfiles}

\begin{document}

% SEGMENT OLP-0246-S01

\olfileid{cmp}{thy}{k1}
\olsection{குறைத்தன்மைக்கு ஓர் எடுத்துக்காட்டு}

\olref[ppr]{prop:reduce} இன் ஒரு பயன்பாட்டைக் கருதுவோம்.

\begin{prop}
\ollabel{prop:k1}
பின்வருமாறு எடுத்துக்கொள்க:
\[
K_1 = \Setabs{e}{\cfind{e}(0) \fdefined}.
\]
அப்போது $K_1$ கணிக்கத்தக்கவாறு எண்ணிடத்தக்கது; ஆனால் கணிக்கத்தக்கது அல்ல.
\end{prop}

\begin{proof}
$K_1 = \Setabs{e}{\lexists[s][T(e,0,s)]}$ என்பதால்,
\olref[eqc]{thm:exists-char} ஆல் $K_1$ கணிக்கத்தக்கவாறு எண்ணிடத்தக்கது.

$K_1$ கணிக்கத்தக்கது அல்ல என்று காட்ட, $K_0$ அதற்குக்
குறைக்கத்தக்கது என்று காட்டுவோம்.

\begin{explain}
இது சற்றுச் சிக்கலானது; ஏனெனில் $K_1$ ஐப் பயன்படுத்தி ஒரு குறிப்பிட்ட
உள்ளீடான $0$ இலிருந்து தொடங்கும் கணிப்புகளைப் பற்றிய கேள்விகளை
மட்டுமே கேட்க முடியும். இவ்வகைக் கேள்விகளுக்கு விடையளிக்கக்கூடிய
அறிவார்ந்த நண்பர் உங்களிடம் இருக்கிறார் எனக் கொள்க (இத்தகைய நண்பர்கள்
``ஆரக்கிள்கள்'' எனப்படுகின்றனர்). இப்போது ஒருவர் உங்களிடம் வந்து,
$\tuple{e, x}$ ஆனது $K_0$ இல் உள்ளதா இல்லையா, அதாவது பொறி $e$
உள்ளீடு $x$ இல் நிற்கிறதா இல்லையா என்று கேட்கிறார் எனக் கொள்க.
நீங்கள் செய்யக்கூடிய ஒன்று, \emph{எந்த} உள்ளீடு வந்தாலும் அதை
ஒதுக்கிவிட்டு, அதற்குப் பதிலாக உள்ளீடு $x$ இல்~$e$ ஐ இயக்கும்
மற்றொரு பொறி~$e_x$ ஐக் கட்டுவது. பொறி $e$ உள்ளீடு $x$ இல்
நிற்கிறதா என்ற கேள்வியும், பொறி $e_x$ உள்ளீடு $0$ இல்
(அல்லது வேறு எந்த உள்ளீட்டிலும்) நிற்கிறதா என்ற கேள்வியும்
தெளிவாகச் சமானமானவை. எனவே இப்புதிய பொறி $e_x$ உள்ளீடு $0$ இல்
நிற்கிறதா என்று உங்கள் நண்பரிடம் கேட்கிறீர்கள்; மாற்றப்பட்ட
கேள்விக்கான நண்பரின் விடை மூலக் கேள்விக்கான விடையைத் தருகிறது.
இது $K_0$ இலிருந்து~$K_1$ க்கு வேண்டிய குறைத்தலை வழங்குகிறது.
\end{explain}

அனைத்தளாவிய பகுதி கணிக்கத்தக்க சார்பைப் பயன்படுத்தி,
$f$ ஐப் பின்வருமாறு வரையறுக்கப்பட்ட 3-இடச் சார்பாக எடுத்துக்கொள்க:
\[
f(x,y,z) \simeq \cfind{x}(y).
\]
$f$ தனது மூன்றாவது உள்ளீட்டை முற்றிலும் ஒதுக்குகிறது என்பதைக்
கவனிக்க. $f = \cfind{e}[3]$ ஆகுமாறு சுட்டெண் $e$ ஒன்றைத் தேர்க;
அப்போது
\[
\cfind{e}[3](x,y,z) \simeq \cfind{x}(y).
\]
$s$-$m$-$n$ தேற்றத்தால், ஒவ்வொரு $z$ க்கும் பின்வருமாறு அமையும்
சார்பு $s(e,x,y)$ ஒன்று உள்ளது:
\begin{align*}
\cfind{s(e,x,y)}(z) & \simeq \cfind{e}[3](x,y,z) \\
& \simeq \cfind{x}(y).
\end{align*}

\begin{explain}
மேலுள்ள முறைசாராத வாதத்தின் மொழியில், $s(e,x,y)$ என்பது எந்த
உள்ளீடு $z$ வந்தாலும் அதை ஒதுக்கிவிட்டு $\cfind{x}(y)$ ஐக்
கணிக்கும் பொறிக்கான ஒரு சுட்டெண்.
\end{explain}

குறிப்பாக,
\[
\cfind{s(e,x,y)}(0) \fdefined \quad \text{எனில், எனில் மட்டுமே} \quad
\cfind{x}(y) \fdefined.
\]
வேறுவிதமாகக் கூறினால், $\tuple{x, y} \in K_0$ எனில், எனில்
மட்டுமே, $s(e,x,y) \in K_1$. எனவே பின்வருமாறு வரையறுக்கப்பட்ட
சார்பு $g$,
\[
g(w) = s(e,(w)_0,(w)_1)
\]
$K_0$ இலிருந்து~$K_1$ க்கான ஒரு குறைத்தல்.
\end{proof}

\end{document}
